\documentclass[11pt]{article}
\usepackage{manual_docs_style}
\hypersetup{hidelinks}
\externaldocument{build/website}[website.pdf]
\externaldocument{build/agentic_rollout}[agentic_rollout.pdf]
\externaldocument{build/dataset_manifest_schema}[dataset_manifest_schema.pdf]

\begin{document}

\docTitle{Public Repository}
\phantomsection\label{docstart:public_repository}
\label{sec:public-repository}

This document describes the public GitHub repository for \name{}.
It's located at \href{https://github.com/TommasoBendinelli/TraceBench.git}{\code{https://github.com/TommasoBendinelli/TraceBench.git}}


\phantomsection\label{sec:public-repository-design-brief}
\section*{Design Brief}

\phantomsection\label{sec:public-repository-overview}
\subsection*{Overview}

It's used to allow user to run rollouts with the same setup as us so that people can be update their results.
It also provides tools to inspect the trajectories and rollouts.
Specifically, the repository should serve four main audiences:

\begin{itemize}
\item who want to download tasks, run an agent, and understand the expected output format.
\item who want to package trajectories and scores for review.
\item people that wants to plot the results.    
\item internal utilities used by GitHub actions.
\end{itemize}

The repository should make these jobs possible without requiring readers to search across private notes, local build artifacts, or the website implementation.
The first page of the repository should answer: what \name{} is, what is downloadable, how to run a minimal workflow, how to submit results, and where to cite the work.

\phantomsection\label{sec:public-repository-positioning}
\subsection*{Positioning}

The public repository should be positioned as a stable benchmark interface, not only as an internal research snapshot.
Its documentation should distinguish the following roles:

\begin{itemize}
\item \textbf{Repository.} Source code, schemas, validators, small examples, scripts, documentation, and release pointers.
\item \textbf{Hugging Face.} Canonical host for large released data, result tables, benchmark bundles, and versioned downloads. Trajectories are hosted there only when included in a published submission archive.
\item \textbf{Website.} Public browsing layer for benchmark explanation, leaderboard inspection, environment pages, and submission guidance.
\item \textbf{Paper.} Methodological and experimental reference.
\end{itemize}

The repository should link to all of these surfaces from the top-level \code{README.md}.
It should avoid duplicating large data files that are already published as release artifacts.

\phantomsection\label{sec:public-repository-public-promise}

\section*{Technical Brief}

\phantomsection\label{sec:public-repository-layout}
\subsection*{Repository layout}

The public repository should use a shallow, readable top-level layout.
The exact names may evolve, but a release-ready repository should include the following categories:

\begin{itemize}
\item \code{README.md}: primary entrypoint for humans.
\item \code{LICENSE}: license terms for code and, where applicable, documentation.
\item \code{CITATION.cff}: machine-readable citation metadata.
\item \code{docs/}: rendered or source documentation for public methods, schemas, and workflows.
\item \code{schemas/}: JSON schemas or schema references for datasets, rollout outputs, submissions, and result records.
\item \code{examples/}: lightweight examples that can be run or inspected without downloading the full benchmark.
\item \code{scripts/} or \code{tools/}: public commands for validation, artifact inspection, and supported downloads.
\item \code{submissions/}: templates and documentation for external result submissions.
\item \code{.github/workflows/}: continuous validation for pull requests and accepted submissions.
\end{itemize}

The repository should not require large binary artifacts to be committed.
When example data is needed, keep it small and clearly separate it from full benchmark releases.

\phantomsection\label{sec:public-repository-readme}
\subsection*{README.md}

The top-level \code{README.md} is the front door of the public repository.
It should be concise, but complete enough that a new user can decide what to do next.
It should include:

\begin{itemize}
\item the benchmark name and one-paragraph summary;
\item links to the paper, website, Hugging Face collection, and documentation;
\item a compact repository map;
\item installation or environment setup instructions for supported public tools;
\item commands for downloading or locating released artifacts;
\item a minimal example workflow;
\item submission instructions and the expected pull-request path;
\item citation, license, contact, and issue-reporting information.
\end{itemize}

The README should distinguish stable commands from development-only commands.
If a command depends on an optional artifact download, the README should say so before the command appears.

\phantomsection\label{sec:public-repository-artifacts}
\subsection*{Artifacts and downloads}

The repository should point to released artifacts rather than duplicating them.
Public release artifacts should be versioned and discoverable from the README and from a dedicated documentation page.
The expected external artifacts are:

\begin{itemize}
\item benchmark task bundles;
\item simulator or environment metadata needed by public examples;
\item full leaderboard result tables;
\item accepted trajectory bundles, when the submitter has released the raw trajectories;
\item validation reports for released submissions;
\item checksums or manifest files for downloaded artifacts.
\end{itemize}

Each release should define which artifact location is canonical.
For \name{}, the default is that Hugging Face stores large artifacts and GitHub stores code, documentation, schemas, validators, and small examples.

\phantomsection\label{sec:public-repository-submissions}
\subsection*{External submissions}

External result submissions should be pull-request based.
The repository should provide a template bundle and validation command that match the rollout and website submission expectations.
Submission documentation should point to the \hyperref[docstart:agentic_rollout]{Agentic Rollout} stage for the trajectory format and to the website technical brief for website-ready records.

A public submission should include:

\begin{itemize}
\item model or agent identity;
\item benchmark version;
\item task or environment coverage;
\item scored outputs;
\item trajectory or artifact pointers;
\item validation metadata;
\item submitter contact and any required declaration.
\end{itemize}

Accepted submissions should be reviewed one submission at a time.
After acceptance, maintainers should publish the canonical accepted artifacts to Hugging Face, regenerate website-ready JSON, and update the public website through the deployment workflow.

\phantomsection\label{sec:public-repository-validation}
\subsection*{Validation and automation}

The public repository should include automated checks for the public interface.
At minimum, GitHub Actions should validate:

\begin{itemize}
\item schema conformance for submitted bundles;
\item required metadata fields;
\item consistency between declared benchmark version and submitted artifacts;
\item basic link or manifest integrity for public release pointers;
\item formatting or lint checks for files that are expected to be machine-read.
\end{itemize}

Validation should fail closed for malformed submissions.
Human maintainers remain responsible for accepting results, publishing large artifacts to Hugging Face, and triggering website updates.

\phantomsection\label{sec:public-repository-boundaries}
\subsection*{Public/private boundaries}

The public repository should not expose local-only, credentialed, or unpublished research state.
The following should remain outside the public repository unless they are intentionally cleaned and documented:

\begin{itemize}
\item private API keys, credentials, tokens, and machine-specific paths;
\item large generated runs that belong on Hugging Face;
\item unpublished or unstable experiments that are not part of a supported release;
\item private review notes and unaccepted third-party submissions;
\item local caches, build directories, and temporary generated files.
\end{itemize}

When an internal file is needed to explain the public workflow, the repository should publish a sanitized schema, template, or example instead of the private source artifact.

\end{document}
